\section{Introduction}
\label{sec:introduction}

Machine learning has been central to high-energy physics for decades, enabling
reconstruction algorithms, event classification, and detector optimization across
collider experiments. A new frontier has emerged with large language models (LLMs)
developed by major AI laboratories, including
Anthropic~\cite{Anthropic:2025accelerating}, OpenAI~\cite{OpenAI:2025gpt5science},
Google DeepMind~\cite{GoogleDeepMind:2025geminiscience}, and others, that can now
assist directly with scientific reasoning, mathematical proof, and code generation.
A striking asymmetry persists: AI is transforming the theoretical and computational side of science research, assisting derivations of new results in theoretical particle physics~\cite{OpenAI:2025newresult}, designing novel protein structures for biology and health research~\cite{GoogleDeepMind:2025alphaproteo}, and simplifying complex calculations. Yet the experimental side of the AI-assisted research loop remains largely absent, especially in fundamental physics. To realize a future where AI rapidly accelerates the discovery of new physics, AI-assisted scientific workflows must be developed and validated across both theory and experiment.

A logical first step toward closing this gap is an agentic analysis of real
experimental data with expert physicist oversight. This is precisely the approach
pursued here. We communicated via natural language prompts with the Anthropic Claude and OpenAI Codex agents, brainstormed
analysis goals, then iteratively prompted the AI agents to execute the full
analysis workflow. No analysis code was written by hand. By analogy, the physicist's role is
similar to advising a graduate student, providing expert direction, reviewing
results, and iterating. The agent was supplied with context comprising the published ALEPH thrust paper~\cite{Heister:2003aj}, prior ROOT-based ALEPH analysis macros covering event selection and thrust calculation, internal analysis notes documenting the miniDST data format and detector calibration procedures, and the archived ALEPH data~\cite{ALEPH:docs}. This prompting strategy follows from our understanding in the AI literature that few-shot examples and supplied context improve AI-assisted outputs~\cite{Brown:2020fewshot,Wei:2022cot}. A detailed inventory of the inputs and prompting approach is given in Section~\ref{sec:ai_agent}.

This vision parallels the evolution of AI in neighboring scientific domains. In
materials science, chemistry, and biology, AI work has progressed from computational
prediction with post-hoc human validation toward AI-designed experiments executed
by autonomous robotic systems, creating closed feedback loops where AI proposes and orchestrates experiments and learns from them. High-energy physics faces a parallel opportunity, with exciting early progress for theory calculations~\cite{Abolhasani:2023sdl,MacLeod:2020sdl,Chen:2024scienceagentbench,Lu:2024aiscientist,Woodruff:2026practicalsi,Bubeck:2025coscientist,OpenAITheory:2026singleminus,Schwartz:2026cparameterllm}.
While we cannot deploy robots to construct particle experiments (yet), the emergence of open and curated data can serve as a
virtual laboratory for developing and validating AI-assisted scientific workflows
that operate across the precision theory--experiment
interface~\cite{Badea:2019vey,Chen:2021uws}.  This requires infrastructure on both sides: precise theoretical predictions and experimental measurements with rigorous uncertainty quantification, where AI-assisted reasoning is tested and held accountable to measurable constraints~\cite{Wang:2023aai,Messeri:2024aiu,Woodruff:2026practicalsi}.

Precision QCD is a natural domain for this program, because $e^{+}e^{-}$ annihilation offers both high-precision data and theory calculations. This matters because fine details are consequential: percent-level uncertainties, coupled to rigorous covariance estimates and end-to-end reproducibility, leave little room for an error to pass unnoticed. In hadronic or nuclear collision environments, by contrast, larger uncertainties can more easily conceal mistakes made by an AI system.

From both an AI and precision physics perspective, thrust is an ideal first observable for this program.
Thrust probes both perturbative radiation and non-perturbative dynamics in a
single distribution, is computable at high perturbative order, and has well-established
connections to resummed calculations, making comparison to theory controlled and
interpretable~\cite{PhysRevLett.39.1587,Abbate:2010xh,Heister:2003aj}. Thrust, along with other event-shape observables measured at LEP, continues to anchor comparisons between
perturbative calculations, generator modeling, and non-perturbative hadronization
constraints. Recent thrust-based $\alpha_s$ extractions report
$\alpha_s(m_Z)\simeq 0.1136\pm 0.0012$, below the 2023 world-average value
$\alpha_s(m_Z)=0.1180\pm 0.0009$~\cite{Benitez:2024nav,ParticleDataGroup:2022pth}.
This persistent tension, often referred to as the $e^{+}e^{-}$ $\alpha_{s}$ puzzle, is of great interest because $\alpha_{s}$ is the least precisely measured coupling constant of the Standard Model. The associated uncertainties propagate to all measurements involving QCD, and the effective field theory tools used to extract $\alpha_{s}$ are widely employed throughout the field. Put differently, if something is missed in the simplest case of $e^{+}e^{-}$ collisions, it is likely that similar effects will also be missed in precision studies of more complex systems. Therefore, testing the new AI technology on thrust in $e^+e^-$ collisions opens the door to significant progress in precision physics.

The objective of the agent-physicist collaboration in this work was to measure the thrust spectrum by applying appropriate selections, unfolding, corrections, and uncertainty handling. The agent was asked to execute the full analysis loop: data and simulation processing, unfolding orchestration,
systematic propagation, figure production, and note maintenance with traceable links
to code and outputs. We utilize the thrust measurement as a testing ground to understand the capabilities of the agents, establish a software base to conduct agentic analysis of the LEP data, and identify areas for improvement. We note that areas for improvement include both the core capabilities of the AI and the interaction of the physicist with the AI.
We envision using our work to pursue a next stage, where the same framework
will be extended to additional event-shape observables (C-parameter and heavy-jet
mass), and the approach of using AI-assisted scientific workflows for both theoretical
and experimental analysis will be pursued in a loop between them. 

The remainder of the document is structured as follows. Section~\ref{sec:ai_agent} defines the AI-agent operating model used in this work.
Section~\ref{sec:aleph} summarizes the ALEPH detector. Section~\ref{sec:datasets} describes data and simulation samples.
Section~\ref{sec:event_selection} details reconstruction and selections, including
particle-category and event-level validation. Section~\ref{sec:method} presents
thrust construction and the executable unfolding workflow.
Section~\ref{sec:systematics} discusses uncertainty measurements. Section~\ref{sec:results} presents the results of the measurement. Section~\ref{sec:conclusion} reports the conclusions from the study. 
%Section~\ref{sec:challenges} closes this work with persistent analysis challenges and mitigation tracking.
